\begin{resumo}
    A organização de partidas de futebol amador em Brasília enfrenta desafios recorrentes, especialmente nas quadras públicas distribuídas pela cidade. A marcação dos jogos costuma depender de grupos dispersos em aplicativos de mensagens, com confirmações inconsistentes e pouca clareza sobre a disponibilidade dos espaços, resultando em cancelamentos, desencontros e subutilização das quadras públicas. Diante dessa problemática, este trabalho propõe o desenvolvimento do BORAFUT, um aplicativo móvel multiplataforma que facilita a organização do futebol amador no Distrito Federal, conectando jogadores, quadras e grupos locais por meio de funcionalidades intuitivas e acessíveis. O desenvolvimento foi fundamentado em uma abordagem metodológica híbrida, combinando elementos de processos dirigidos a plano e metodologias ágeis, com base nos princípios do Extreme Programming (XP). O processo de levantamento de requisitos foi realizado por meio dos métodos de observação, \textit{brainwriting} e um questionário aplicado a mais de 60 praticantes de futebol amador, e o processo de priorização e modelagem foi realizado utilizando o método MoSCoW e a decomposição hierárquica do backlog em níveis sucessivos de granularidade, dos temas estratégicos às tarefas técnicas. O objetivo deste trabalho é abranger de forma integrada o ciclo completo de desenvolvimento de software no contexto de um produto mínimo viável, desde a investigação do problema, definição e modelagem dos requisitos até a concepção da arquitetura, do banco de dados e da identidade visual, resultando na produção de um protótipo interativo de alta fidelidade que funciona como prova de conceito das principais jornadas do BORAFUT e fornecendo uma base sólida para as etapas de implementação e validação da solução proposta.

    \vspace{\onelineskip}
    
    \noindent
    \textbf{Palavras-chave}: futebol amador. aplicativo móvel. engenharia de software. metodologias ágeis. espaços públicos.
\end{resumo}
